% Options for packages loaded elsewhere
\PassOptionsToPackage{unicode}{hyperref}
\PassOptionsToPackage{hyphens}{url}
%
\documentclass[
  ignorenonframetext,
  aspectratio=169,
]{beamer}
\usepackage{pgfpages}
\setbeamertemplate{caption}[numbered]
\setbeamertemplate{caption label separator}{: }
\setbeamercolor{caption name}{fg=normal text.fg}
\beamertemplatenavigationsymbolsempty
% Prevent slide breaks in the middle of a paragraph
\widowpenalties 1 10000
\raggedbottom
\setbeamertemplate{part page}{
  \centering
  \begin{beamercolorbox}[sep=16pt,center]{part title}
    \usebeamerfont{part title}\insertpart\par
  \end{beamercolorbox}
}
\setbeamertemplate{section page}{
  \centering
  \begin{beamercolorbox}[sep=12pt,center]{part title}
    \usebeamerfont{section title}\insertsection\par
  \end{beamercolorbox}
}
\setbeamertemplate{subsection page}{
  \centering
  \begin{beamercolorbox}[sep=8pt,center]{part title}
    \usebeamerfont{subsection title}\insertsubsection\par
  \end{beamercolorbox}
}
\AtBeginPart{
  \frame{\partpage}
}
\AtBeginSection{
  \ifbibliography
  \else
    \frame{\sectionpage}
  \fi
}
\AtBeginSubsection{
  \frame{\subsectionpage}
}

\usepackage{amsmath,amssymb}
\usepackage{iftex}
\ifPDFTeX
  \usepackage[T1]{fontenc}
  \usepackage[utf8]{inputenc}
  \usepackage{textcomp} % provide euro and other symbols
\else % if luatex or xetex
  \usepackage{unicode-math}
  \defaultfontfeatures{Scale=MatchLowercase}
  \defaultfontfeatures[\rmfamily]{Ligatures=TeX,Scale=1}
\fi
\usepackage{lmodern}
\usetheme[]{Madrid}
\usecolortheme{default}
\ifPDFTeX\else  
    % xetex/luatex font selection
\fi
% Use upquote if available, for straight quotes in verbatim environments
\IfFileExists{upquote.sty}{\usepackage{upquote}}{}
\IfFileExists{microtype.sty}{% use microtype if available
  \usepackage[]{microtype}
  \UseMicrotypeSet[protrusion]{basicmath} % disable protrusion for tt fonts
}{}
\makeatletter
\@ifundefined{KOMAClassName}{% if non-KOMA class
  \IfFileExists{parskip.sty}{%
    \usepackage{parskip}
  }{% else
    \setlength{\parindent}{0pt}
    \setlength{\parskip}{6pt plus 2pt minus 1pt}}
}{% if KOMA class
  \KOMAoptions{parskip=half}}
\makeatother
\usepackage{xcolor}
\newif\ifbibliography
\setlength{\emergencystretch}{3em} % prevent overfull lines
\setcounter{secnumdepth}{5}

\usepackage{color}
\usepackage{fancyvrb}
\newcommand{\VerbBar}{|}
\newcommand{\VERB}{\Verb[commandchars=\\\{\}]}
\DefineVerbatimEnvironment{Highlighting}{Verbatim}{commandchars=\\\{\}}
% Add ',fontsize=\small' for more characters per line
\usepackage{framed}
\definecolor{shadecolor}{RGB}{241,243,245}
\newenvironment{Shaded}{\begin{snugshade}}{\end{snugshade}}
\newcommand{\AlertTok}[1]{\textcolor[rgb]{0.68,0.00,0.00}{#1}}
\newcommand{\AnnotationTok}[1]{\textcolor[rgb]{0.37,0.37,0.37}{#1}}
\newcommand{\AttributeTok}[1]{\textcolor[rgb]{0.40,0.45,0.13}{#1}}
\newcommand{\BaseNTok}[1]{\textcolor[rgb]{0.68,0.00,0.00}{#1}}
\newcommand{\BuiltInTok}[1]{\textcolor[rgb]{0.00,0.23,0.31}{#1}}
\newcommand{\CharTok}[1]{\textcolor[rgb]{0.13,0.47,0.30}{#1}}
\newcommand{\CommentTok}[1]{\textcolor[rgb]{0.37,0.37,0.37}{#1}}
\newcommand{\CommentVarTok}[1]{\textcolor[rgb]{0.37,0.37,0.37}{\textit{#1}}}
\newcommand{\ConstantTok}[1]{\textcolor[rgb]{0.56,0.35,0.01}{#1}}
\newcommand{\ControlFlowTok}[1]{\textcolor[rgb]{0.00,0.23,0.31}{#1}}
\newcommand{\DataTypeTok}[1]{\textcolor[rgb]{0.68,0.00,0.00}{#1}}
\newcommand{\DecValTok}[1]{\textcolor[rgb]{0.68,0.00,0.00}{#1}}
\newcommand{\DocumentationTok}[1]{\textcolor[rgb]{0.37,0.37,0.37}{\textit{#1}}}
\newcommand{\ErrorTok}[1]{\textcolor[rgb]{0.68,0.00,0.00}{#1}}
\newcommand{\ExtensionTok}[1]{\textcolor[rgb]{0.00,0.23,0.31}{#1}}
\newcommand{\FloatTok}[1]{\textcolor[rgb]{0.68,0.00,0.00}{#1}}
\newcommand{\FunctionTok}[1]{\textcolor[rgb]{0.28,0.35,0.67}{#1}}
\newcommand{\ImportTok}[1]{\textcolor[rgb]{0.00,0.46,0.62}{#1}}
\newcommand{\InformationTok}[1]{\textcolor[rgb]{0.37,0.37,0.37}{#1}}
\newcommand{\KeywordTok}[1]{\textcolor[rgb]{0.00,0.23,0.31}{#1}}
\newcommand{\NormalTok}[1]{\textcolor[rgb]{0.00,0.23,0.31}{#1}}
\newcommand{\OperatorTok}[1]{\textcolor[rgb]{0.37,0.37,0.37}{#1}}
\newcommand{\OtherTok}[1]{\textcolor[rgb]{0.00,0.23,0.31}{#1}}
\newcommand{\PreprocessorTok}[1]{\textcolor[rgb]{0.68,0.00,0.00}{#1}}
\newcommand{\RegionMarkerTok}[1]{\textcolor[rgb]{0.00,0.23,0.31}{#1}}
\newcommand{\SpecialCharTok}[1]{\textcolor[rgb]{0.37,0.37,0.37}{#1}}
\newcommand{\SpecialStringTok}[1]{\textcolor[rgb]{0.13,0.47,0.30}{#1}}
\newcommand{\StringTok}[1]{\textcolor[rgb]{0.13,0.47,0.30}{#1}}
\newcommand{\VariableTok}[1]{\textcolor[rgb]{0.07,0.07,0.07}{#1}}
\newcommand{\VerbatimStringTok}[1]{\textcolor[rgb]{0.13,0.47,0.30}{#1}}
\newcommand{\WarningTok}[1]{\textcolor[rgb]{0.37,0.37,0.37}{\textit{#1}}}

\providecommand{\tightlist}{%
  \setlength{\itemsep}{0pt}\setlength{\parskip}{0pt}}\usepackage{longtable,booktabs,array}
\usepackage{calc} % for calculating minipage widths
\usepackage{caption}
% Make caption package work with longtable
\makeatletter
\def\fnum@table{\tablename~\thetable}
\makeatother
\usepackage{graphicx}
\makeatletter
\def\maxwidth{\ifdim\Gin@nat@width>\linewidth\linewidth\else\Gin@nat@width\fi}
\def\maxheight{\ifdim\Gin@nat@height>\textheight\textheight\else\Gin@nat@height\fi}
\makeatother
% Scale images if necessary, so that they will not overflow the page
% margins by default, and it is still possible to overwrite the defaults
% using explicit options in \includegraphics[width, height, ...]{}
\setkeys{Gin}{width=\maxwidth,height=\maxheight,keepaspectratio}
% Set default figure placement to htbp
\makeatletter
\def\fps@figure{htbp}
\makeatother

\usepackage{etoolbox}
\AtBeginEnvironment{Shaded}{\scriptsize}
\makeatletter
\def\verbatim@font{\scriptsize\ttfamily}
\makeatother
\AtBeginSection[]{\begin{frame}{Outline}\tableofcontents[currentsection]\end{frame}}
\makeatletter
\makeatother
\makeatletter
\makeatother
\makeatletter
\@ifpackageloaded{caption}{}{\usepackage{caption}}
\AtBeginDocument{%
\ifdefined\contentsname
  \renewcommand*\contentsname{Table of contents}
\else
  \newcommand\contentsname{Table of contents}
\fi
\ifdefined\listfigurename
  \renewcommand*\listfigurename{List of Figures}
\else
  \newcommand\listfigurename{List of Figures}
\fi
\ifdefined\listtablename
  \renewcommand*\listtablename{List of Tables}
\else
  \newcommand\listtablename{List of Tables}
\fi
\ifdefined\figurename
  \renewcommand*\figurename{Figure}
\else
  \newcommand\figurename{Figure}
\fi
\ifdefined\tablename
  \renewcommand*\tablename{Table}
\else
  \newcommand\tablename{Table}
\fi
}
\@ifpackageloaded{float}{}{\usepackage{float}}
\floatstyle{ruled}
\@ifundefined{c@chapter}{\newfloat{codelisting}{h}{lop}}{\newfloat{codelisting}{h}{lop}[chapter]}
\floatname{codelisting}{Listing}
\newcommand*\listoflistings{\listof{codelisting}{List of Listings}}
\makeatother
\makeatletter
\@ifpackageloaded{caption}{}{\usepackage{caption}}
\@ifpackageloaded{subcaption}{}{\usepackage{subcaption}}
\makeatother
\makeatletter
\@ifpackageloaded{tcolorbox}{}{\usepackage[skins,breakable]{tcolorbox}}
\makeatother
\makeatletter
\@ifundefined{shadecolor}{\definecolor{shadecolor}{rgb}{.97, .97, .97}}
\makeatother
\makeatletter
\makeatother
\makeatletter
\makeatother
\ifLuaTeX
  \usepackage{selnolig}  % disable illegal ligatures
\fi
\IfFileExists{bookmark.sty}{\usepackage{bookmark}}{\usepackage{hyperref}}
\IfFileExists{xurl.sty}{\usepackage{xurl}}{} % add URL line breaks if available
\urlstyle{same} % disable monospaced font for URLs
\hypersetup{
  pdftitle={Session 1: Setting Up R and RStudio},
  pdfauthor={Taikgun Song},
  hidelinks,
  pdfcreator={LaTeX via pandoc}}

\title{Session 1: Setting Up R and RStudio}
\subtitle{R Workshop for DBA Students}
\author{Taikgun Song}
\date{2026-11-06}

\begin{document}
\frame{\titlepage}
\ifdefined\Shaded\renewenvironment{Shaded}{\begin{tcolorbox}[frame hidden, sharp corners, breakable, borderline west={3pt}{0pt}{shadecolor}, boxrule=0pt, enhanced, interior hidden]}{\end{tcolorbox}}\fi

\renewcommand*\contentsname{Outline}
\begin{frame}[allowframebreaks]
  \frametitle{Outline}
  \tableofcontents[hideallsubsections]
\end{frame}
\begin{frame}
\end{frame}

\begin{frame}{Session Objectives}
\protect\hypertarget{session-objectives}{}
By the end of this session, you will be able to:

\begin{enumerate}
\tightlist
\item
  Install the correct versions of R and RStudio for your computer
\item
  Configure RStudio so that your work is reproducible from the start
\item
  Install, load, and update R packages
\item
  Explain working directories and file paths, and organize work with
  RStudio Projects
\item
  Create, find, and read files from R
\item
  Render your first Quarto document
\item
  Find help on your own when something goes wrong
\end{enumerate}
\end{frame}

\hypertarget{sec-what}{%
\section{R and RStudio}\label{sec-what}}

\begin{frame}{R and RStudio}
\protect\hypertarget{r-and-rstudio}{}
\begin{columns}[T]
\begin{column}{0.5\textwidth}
\textbf{R}

\begin{itemize}
\tightlist
\item
  A programming language for statistics and data analysis
\item
  Free and open source
\item
  Most new statistical methods appear first as R packages
\end{itemize}
\end{column}

\begin{column}{0.5\textwidth}
\textbf{RStudio}

\begin{itemize}
\tightlist
\item
  A program that makes R easier to use
\item
  Script editor, file browser, plots, help: one window
\item
  Made by a company called Posit
\end{itemize}
\end{column}
\end{columns}

\begin{exampleblock}{R is the engine, RStudio is the dashboard}
\protect\hypertarget{r-is-the-engine-rstudio-is-the-dashboard}{}
You need both, and you install them separately. You open RStudio, and
RStudio runs R for you.
\end{exampleblock}
\end{frame}

\hypertarget{sec-install}{%
\section{Installing R and RStudio}\label{sec-install}}

\begin{frame}[fragile]{Step 1: Know Your Computer}
\protect\hypertarget{sec-know-computer}{}
\begin{longtable}[]{@{}
  >{\raggedright\arraybackslash}p{(\columnwidth - 4\tabcolsep) * \real{0.2439}}
  >{\raggedright\arraybackslash}p{(\columnwidth - 4\tabcolsep) * \real{0.3415}}
  >{\raggedright\arraybackslash}p{(\columnwidth - 4\tabcolsep) * \real{0.4146}}@{}}
\toprule\noalign{}
\begin{minipage}[b]{\linewidth}\raggedright
Computer
\end{minipage} & \begin{minipage}[b]{\linewidth}\raggedright
How to check
\end{minipage} & \begin{minipage}[b]{\linewidth}\raggedright
What to install
\end{minipage} \\
\midrule\noalign{}
\endhead
\textbf{Windows} & Settings \(\rightarrow\) System \(\rightarrow\) About
\(\rightarrow\) System type & \texttt{x64} for almost everyone \\
\textbf{Mac} & Apple menu \(\rightarrow\) About This Mac & \textbf{Chip:
Apple M\ldots{}} \(\rightarrow\) \texttt{arm64} \newline **Processor:
Intel** \(\rightarrow\) \texttt{x86\_64} \\
\bottomrule\noalign{}
\end{longtable}

\begin{alertblock}{Mac users}
\protect\hypertarget{mac-users}{}
\begin{itemize}
\tightlist
\item
  The \texttt{arm64} installer needs \textbf{macOS 14 (Sonoma) or later}
  (R 4.6.2, Nov 2026)
\item
  Do \textbf{not} install the Intel version on an Apple silicon Mac: it
  runs through Rosetta, is slower, and some packages fail in confusing
  ways
\item
  After installing, check in the Console: \texttt{R.version\$arch} must
  say \texttt{"aarch64"} on Apple silicon (\texttt{"x86\_64"} = Intel)
\end{itemize}
\end{alertblock}
\end{frame}

\begin{frame}[fragile]{Steps 2 and 3: Install R, then RStudio}
\protect\hypertarget{steps-2-and-3-install-r-then-rstudio}{}
\textbf{R} from CRAN: \url{https://cloud.r-project.org/}

\begin{itemize}
\tightlist
\item
  \textbf{Windows:} Download R for Windows \(\rightarrow\) base
  \(\rightarrow\) run the installer, accept defaults
\item
  \textbf{Mac:} Download R for macOS \(\rightarrow\) the \texttt{.pkg}
  that matches your chip
\end{itemize}

\textbf{RStudio Desktop} (free):
\url{https://posit.co/download/rstudio-desktop/}

\begin{itemize}
\tightlist
\item
  \textbf{Windows:} run the installer
\item
  \textbf{Mac:} open the \texttt{.dmg}, drag RStudio into
  \textbf{Applications}
\end{itemize}

RStudio comes with \textbf{Quarto} built in.

\textbf{Alternative:} \href{https://positron.posit.co/}{Positron}, a
newer editor from Posit for R and Python. This course uses RStudio.
\end{frame}

\begin{frame}{Keep R Up to Date}
\protect\hypertarget{keep-r-up-to-date}{}
\begin{itemize}
\tightlist
\item
  A new version of R comes out every April, with bug-fix releases during
  the year
\item
  Updating = installing the new version over the old one
\item
  After a major update (e.g.~4.5 \(\rightarrow\) 4.6), \textbf{reinstall
  your packages}
\end{itemize}
\end{frame}

\begin{frame}[fragile]{Special Cases}
\protect\hypertarget{special-cases}{}
\textbf{Windows on ARM} (Snapdragon ``Copilot+ PCs'')

\begin{itemize}
\tightlist
\item
  Native R for Windows 11 exists, but is still \textbf{experimental};
  some packages fail
\item
  Most reliable option: \textbf{Posit Cloud}
\end{itemize}

\textbf{Compiling packages from source} (rarely needed)

\begin{itemize}
\tightlist
\item
  Windows: \textbf{Rtools}, matching your R version
\item
  Mac: Command Line Tools: run \texttt{xcode-select\ -\/-install} in
  Terminal
\item
  Install only if an error message tells you to
\end{itemize}

\begin{Shaded}
\begin{Highlighting}[]
\FunctionTok{install.packages}\NormalTok{(}\StringTok{"pkgbuild"}\NormalTok{)              }\CommentTok{\# optional check after installing them}
\NormalTok{pkgbuild}\SpecialCharTok{::}\FunctionTok{has\_build\_tools}\NormalTok{(}\AttributeTok{debug =} \ConstantTok{TRUE}\NormalTok{)   }\CommentTok{\# TRUE = ready to compile packages}
\end{Highlighting}
\end{Shaded}
\end{frame}

\begin{frame}{Posit Cloud: A Backup Option}
\protect\hypertarget{sec-cloud}{}
\url{https://posit.cloud}: RStudio in your web browser, nothing to
install

Useful when:

\begin{itemize}
\tightlist
\item
  Your \textbf{work laptop} does not allow installing software
\item
  You have a Windows ARM laptop or an old operating system
\item
  Your installation breaks and you need to keep working
\end{itemize}

Free plan \(\rightarrow\) \textbf{New Project \(\rightarrow\) New
RStudio Project} (check current usage limits)

\begin{exampleblock}{For your own research}
\protect\hypertarget{for-your-own-research}{}
A local installation is still better: faster, no usage limits, and your
data stays on your computer.
\end{exampleblock}
\end{frame}

\hypertarget{sec-tour}{%
\section{A Tour of RStudio}\label{sec-tour}}

\begin{frame}{The Four Panes}
\protect\hypertarget{the-four-panes}{}
\includegraphics[width=0.75\textwidth,height=\textheight]{images/rstudio_panes.png}
\end{frame}

\begin{frame}[fragile]{What Each Pane Is For}
\protect\hypertarget{what-each-pane-is-for}{}
\begin{longtable}[]{@{}
  >{\raggedright\arraybackslash}p{(\columnwidth - 4\tabcolsep) * \real{0.2069}}
  >{\raggedright\arraybackslash}p{(\columnwidth - 4\tabcolsep) * \real{0.2414}}
  >{\raggedright\arraybackslash}p{(\columnwidth - 4\tabcolsep) * \real{0.5517}}@{}}
\toprule\noalign{}
\begin{minipage}[b]{\linewidth}\raggedright
Pane
\end{minipage} & \begin{minipage}[b]{\linewidth}\raggedright
Where
\end{minipage} & \begin{minipage}[b]{\linewidth}\raggedright
What it is for
\end{minipage} \\
\midrule\noalign{}
\endhead
\textbf{Source} & Top left & Writing and saving scripts (\texttt{.R})
and documents (\texttt{.qmd}) \\
\textbf{Console} & Bottom left & Where R runs code; quick, throwaway
commands \\
\textbf{Environment} & Top right & Objects currently in memory \\
\textbf{Files / Plots / Packages / Help} & Bottom right & Files, plots,
packages, help pages \\
\bottomrule\noalign{}
\end{longtable}

\begin{exampleblock}{Script or Console?}
\protect\hypertarget{script-or-console}{}
Anything you want to keep goes in a \textbf{script} or Quarto document.
If it only exists in the Console, it is gone when you close RStudio.
\end{exampleblock}
\end{frame}

\begin{frame}[fragile]{Keyboard Shortcuts}
\protect\hypertarget{keyboard-shortcuts}{}
\begin{longtable}[]{@{}
  >{\raggedright\arraybackslash}p{(\columnwidth - 4\tabcolsep) * \real{0.3636}}
  >{\raggedright\arraybackslash}p{(\columnwidth - 4\tabcolsep) * \real{0.4091}}
  >{\raggedright\arraybackslash}p{(\columnwidth - 4\tabcolsep) * \real{0.2273}}@{}}
\toprule\noalign{}
\begin{minipage}[b]{\linewidth}\raggedright
Action
\end{minipage} & \begin{minipage}[b]{\linewidth}\raggedright
Windows
\end{minipage} & \begin{minipage}[b]{\linewidth}\raggedright
Mac
\end{minipage} \\
\midrule\noalign{}
\endhead
Run current line or selection & \texttt{Ctrl} + \texttt{Enter} &
\texttt{Cmd} + \texttt{Return} \\
Run the whole script & \texttt{Ctrl} + \texttt{Shift} + \texttt{Enter} &
\texttt{Cmd} + \texttt{Shift} + \texttt{Return} \\
New R script & \texttt{Ctrl} + \texttt{Shift} + \texttt{N} &
\texttt{Cmd} + \texttt{Shift} + \texttt{N} \\
Comment / uncomment & \texttt{Ctrl} + \texttt{Shift} + \texttt{C} &
\texttt{Cmd} + \texttt{Shift} + \texttt{C} \\
Clear the Console & \texttt{Ctrl} + \texttt{L} & \texttt{Ctrl} +
\texttt{L} \\
Stop running code & \texttt{Esc} & \texttt{Esc} \\
Show all shortcuts & \texttt{Alt} + \texttt{Shift} + \texttt{K} &
\texttt{Option} + \texttt{Shift} + \texttt{K} \\
\bottomrule\noalign{}
\end{longtable}
\end{frame}

\begin{frame}{Settings to Change on Day One}
\protect\hypertarget{sec-settings}{}
\textbf{Tools \(\rightarrow\) Global Options}

\begin{longtable}[]{@{}
  >{\raggedright\arraybackslash}p{(\columnwidth - 4\tabcolsep) * \real{0.3333}}
  >{\raggedright\arraybackslash}p{(\columnwidth - 4\tabcolsep) * \real{0.3000}}
  >{\raggedright\arraybackslash}p{(\columnwidth - 4\tabcolsep) * \real{0.3667}}@{}}
\toprule\noalign{}
\begin{minipage}[b]{\linewidth}\raggedright
Location
\end{minipage} & \begin{minipage}[b]{\linewidth}\raggedright
Setting
\end{minipage} & \begin{minipage}[b]{\linewidth}\raggedright
Change to
\end{minipage} \\
\midrule\noalign{}
\endhead
General \(\rightarrow\) Basic & Restore .RData into workspace at startup
& \textbf{Unchecked} \\
General \(\rightarrow\) Basic & Save workspace to .RData on exit &
\textbf{Never} \\
Code \(\rightarrow\) Saving & Default text encoding & \textbf{UTF-8} \\
\bottomrule\noalign{}
\end{longtable}
\end{frame}

\begin{frame}[fragile]{Why Not Save the Workspace?}
\protect\hypertarget{why-not-save-the-workspace}{}
\begin{itemize}
\tightlist
\item
  By default, RStudio saves everything in memory on exit and reloads it
  next time
\item
  Then your results can depend on objects you created weeks ago and
  \textbf{no longer have code for}
\end{itemize}

\begin{alertblock}{Reproducibility}
\protect\hypertarget{reproducibility}{}
Your \textbf{script} should be the complete record of your analysis. If
you can restart R, rerun the script, and get the same results, your work
is reproducible.

\begin{itemize}
\tightlist
\item
  Restart R any time: \textbf{Session \(\rightarrow\) Restart R}
  (\texttt{Ctrl}/\texttt{Cmd} + \texttt{Shift} + \texttt{F10})
\end{itemize}
\end{alertblock}
\end{frame}

\hypertarget{sec-packages}{%
\section{R Packages}\label{sec-packages}}

\begin{frame}[fragile]{What Is a Package?}
\protect\hypertarget{what-is-a-package}{}
\begin{itemize}
\tightlist
\item
  A bundle of functions, data, and documentation that someone has shared
\item
  More than 20,000 on CRAN alone
\item
  Like an app on your phone: \textbf{install} once, \textbf{open (load)}
  each time
\end{itemize}

\begin{longtable}[]{@{}
  >{\raggedright\arraybackslash}p{(\columnwidth - 4\tabcolsep) * \real{0.2222}}
  >{\raggedright\arraybackslash}p{(\columnwidth - 4\tabcolsep) * \real{0.3704}}
  >{\raggedright\arraybackslash}p{(\columnwidth - 4\tabcolsep) * \real{0.4074}}@{}}
\toprule\noalign{}
\begin{minipage}[b]{\linewidth}\raggedright
Step
\end{minipage} & \begin{minipage}[b]{\linewidth}\raggedright
Function
\end{minipage} & \begin{minipage}[b]{\linewidth}\raggedright
How often
\end{minipage} \\
\midrule\noalign{}
\endhead
Install & \texttt{install.packages("name")} & Once per computer (and
after updating R) \\
Load & \texttt{library(name)} & Every session, at the top of every
script \\
\bottomrule\noalign{}
\end{longtable}
\end{frame}

\begin{frame}[fragile]{Installing Packages}
\protect\hypertarget{installing-packages}{}
\begin{Shaded}
\begin{Highlighting}[]
\FunctionTok{install.packages}\NormalTok{(}\StringTok{"readxl"}\NormalTok{)                       }\CommentTok{\# one package (note the quotes)}
\FunctionTok{install.packages}\NormalTok{(}\FunctionTok{c}\NormalTok{(}\StringTok{"readxl"}\NormalTok{, }\StringTok{"haven"}\NormalTok{, }\StringTok{"here"}\NormalTok{))   }\CommentTok{\# several at once}
\end{Highlighting}
\end{Shaded}

Or: \textbf{Packages} pane \(\rightarrow\) \textbf{Install}
\(\rightarrow\) type the name

\begin{exampleblock}{If R asks a question during installation}
\protect\hypertarget{if-r-asks-a-question-during-installation}{}
\begin{itemize}
\tightlist
\item
  \emph{``Would you like to use a personal library instead?''}
  \(\rightarrow\) \textbf{yes} (common on work laptops)
\item
  \emph{``Do you want to install from sources the packages which need
  compilation?''} \(\rightarrow\) \textbf{no}
\end{itemize}
\end{exampleblock}
\end{frame}

\begin{frame}[fragile]{The Course Install List}
\protect\hypertarget{sec-course-packages}{}
\begin{Shaded}
\begin{Highlighting}[]
\FunctionTok{install.packages}\NormalTok{(}\FunctionTok{c}\NormalTok{(}
  \StringTok{"tidyverse"}\NormalTok{,     }\CommentTok{\# data wrangling and visualization}
  \StringTok{"here"}\NormalTok{,          }\CommentTok{\# file paths that work on every computer}
  \StringTok{"readxl"}\NormalTok{,        }\CommentTok{\# read Excel files}
  \StringTok{"haven"}\NormalTok{,         }\CommentTok{\# read SPSS, Stata, and SAS files}
  \StringTok{"car"}\NormalTok{,           }\CommentTok{\# regression tools; also installs carData}
  \StringTok{"gtsummary"}\NormalTok{,     }\CommentTok{\# publication{-}ready summary tables}
  \StringTok{"modelsummary"}\NormalTok{,  }\CommentTok{\# publication{-}ready regression tables}
  \StringTok{"rmarkdown"}\NormalTok{,     }\CommentTok{\# used by Quarto to run R code}
  \StringTok{"quarto"}         \CommentTok{\# control Quarto from R}
\NormalTok{))}
\end{Highlighting}
\end{Shaded}
\end{frame}

\begin{frame}[fragile]{Check the Installation}
\protect\hypertarget{check-the-installation}{}
\begin{Shaded}
\begin{Highlighting}[]
\NormalTok{course\_pkgs }\OtherTok{\textless{}{-}} \FunctionTok{c}\NormalTok{(}\StringTok{"tidyverse"}\NormalTok{, }\StringTok{"here"}\NormalTok{, }\StringTok{"readxl"}\NormalTok{, }\StringTok{"haven"}\NormalTok{, }\StringTok{"car"}\NormalTok{,}
                 \StringTok{"gtsummary"}\NormalTok{, }\StringTok{"modelsummary"}\NormalTok{, }\StringTok{"rmarkdown"}\NormalTok{, }\StringTok{"quarto"}\NormalTok{)}

\FunctionTok{sapply}\NormalTok{(course\_pkgs, requireNamespace, }\AttributeTok{quietly =} \ConstantTok{TRUE}\NormalTok{)   }\CommentTok{\# TRUE = installed}
\end{Highlighting}
\end{Shaded}

\begin{verbatim}
   tidyverse         here       readxl        haven          car 
        TRUE         TRUE         TRUE         TRUE         TRUE 
   gtsummary modelsummary    rmarkdown       quarto 
       FALSE        FALSE         TRUE        FALSE 
\end{verbatim}

\begin{itemize}
\item
  \texttt{\textless{}-} stores the list under the name
  \texttt{course\_pkgs}; Session 3 covers operators in detail
\item
  Every value should be \texttt{TRUE}
\item
  For any \texttt{FALSE}: run \texttt{install.packages()} again and
  \textbf{read the error message}
\end{itemize}
\end{frame}

\begin{frame}[fragile]{Loading Packages}
\protect\hypertarget{loading-packages}{}
\begin{Shaded}
\begin{Highlighting}[]
\FunctionTok{library}\NormalTok{(tidyverse)}
\end{Highlighting}
\end{Shaded}

\begin{verbatim}
Warning: package 'tidyverse' was built under R version 4.3.3
\end{verbatim}

\begin{verbatim}
-- Attaching core tidyverse packages -------------- tidyverse 2.0.0 --
v dplyr     1.1.2     v readr     2.2.0
v forcats   1.0.1     v stringr   1.5.0
v ggplot2   3.4.2     v tibble    3.2.1
v lubridate 1.9.2     v tidyr     1.3.0
v purrr     1.0.1     
-- Conflicts -------------------------------- tidyverse_conflicts() --
x dplyr::filter() masks stats::filter()
x dplyr::lag()    masks stats::lag()
i Use the conflicted package (<http://conflicted.r-lib.org/>) to force all conflicts to become errors
\end{verbatim}

\begin{itemize}
\tightlist
\item
  \textbf{Conflicts}: \texttt{dplyr::filter()} masks
  \texttt{stats::filter()}, so \texttt{filter()} now means the
  \texttt{dplyr} version. This is normal, not an error
\end{itemize}
\end{frame}

\begin{frame}[fragile]{\texttt{package::function()} and
\texttt{require()}}
\protect\hypertarget{packagefunction-and-require}{}
Use one function without loading the whole package:

\begin{Shaded}
\begin{Highlighting}[]
\NormalTok{here}\SpecialCharTok{::}\FunctionTok{here}\NormalTok{()}
\end{Highlighting}
\end{Shaded}

\begin{verbatim}
[1] "C:/Users/USER/Desktop/DBA_R_Workshop"
\end{verbatim}

\begin{itemize}
\tightlist
\item
  \texttt{::} also shows where a function comes from
\end{itemize}

\textbf{\texttt{library()} vs.~\texttt{require()}}

\begin{itemize}
\tightlist
\item
  \texttt{require()} returns \texttt{FALSE} with a warning when a
  package is missing
\item
  \texttt{library()} stops with an error: \textbf{use it in scripts}
\end{itemize}
\end{frame}

\begin{frame}[fragile]{Updating and Removing Packages}
\protect\hypertarget{updating-and-removing-packages}{}
\begin{Shaded}
\begin{Highlighting}[]
\FunctionTok{update.packages}\NormalTok{(}\AttributeTok{ask =} \ConstantTok{FALSE}\NormalTok{)   }\CommentTok{\# update everything (or Packages pane \textgreater{} Update)}
\FunctionTok{remove.packages}\NormalTok{(}\StringTok{"readxl"}\NormalTok{)      }\CommentTok{\# uninstall}
\end{Highlighting}
\end{Shaded}

\begin{itemize}
\tightlist
\item
  \texttt{install.packages()} installs from \textbf{CRAN}, where every
  package is checked
\item
  Packages can also come from other repositories, such as
  \textbf{GitHub} (in development, not checked by CRAN) or
  \textbf{Bioconductor} (bioinformatics), each with its own installation
  method
\end{itemize}
\end{frame}

\begin{frame}[fragile]{\texttt{sessionInfo()}}
\protect\hypertarget{sec-session-info}{}
\begin{Shaded}
\begin{Highlighting}[]
\FunctionTok{sessionInfo}\NormalTok{()      }\CommentTok{\# output is long: run it in the Console to see yours}
\end{Highlighting}
\end{Shaded}

\begin{itemize}
\tightlist
\item
  Your R version, operating system, and loaded packages with versions
\item
  \textbf{Include it whenever you ask for help}
\end{itemize}
\end{frame}

\hypertarget{sec-paths}{%
\section{Working Directories and File Paths}\label{sec-paths}}

\begin{frame}[fragile]{Folders Form a Tree}
\protect\hypertarget{folders-form-a-tree}{}
\begin{verbatim}
C:/                              <- the root: top of the tree (Mac: /)
`-- Users/
    `-- tek/                     <- your home folder (~)
        `-- Documents/
            `-- dba/             <- suppose this is the working directory
                |-- analysis.qmd
                `-- data/
                    `-- survey.csv
\end{verbatim}

\begin{itemize}
\tightlist
\item
  \textbf{Absolute} path, from the root:
  \texttt{C:/Users/tek/Documents/dba/data/survey.csv}
\item
  \textbf{Relative} path, from \texttt{dba/}: \texttt{data/survey.csv}
\item
  Going \textbf{up} with \texttt{..}, from \texttt{data/}:
  \texttt{../analysis.qmd}
\item
  Mac: same tree, but the root is \texttt{/} with no drive letter:
  \texttt{/Users/tek/Documents/dba/...}
\end{itemize}
\end{frame}

\begin{frame}[fragile]{Folders and Paths}
\protect\hypertarget{folders-and-paths}{}
\begin{longtable}[]{@{}ll@{}}
\toprule\noalign{}
Kind of path & Example \\
\midrule\noalign{}
\endhead
Absolute (Windows) &
\texttt{C:/Users/tek/Documents/dba/data/survey.csv} \\
Absolute (Mac) & \texttt{/Users/tek/Documents/dba/data/survey.csv} \\
Relative & \texttt{data/survey.csv} (from the \textbf{current
folder}) \\
\bottomrule\noalign{}
\end{longtable}

\begin{longtable}[]{@{}ll@{}}
\toprule\noalign{}
Symbol & Meaning \\
\midrule\noalign{}
\endhead
\texttt{/} & Separates folders; the root folder on a Mac \\
\texttt{\textasciitilde{}} & Your home folder \\
\texttt{.} / \texttt{..} & The current folder / one level up \\
\bottomrule\noalign{}
\end{longtable}

\begin{alertblock}{Windows users}
\protect\hypertarget{windows-users}{}
Use forward slashes (\texttt{C:/Users/tek}) or double backslashes
(\texttt{C:\textbackslash{}\textbackslash{}Users\textbackslash{}\textbackslash{}tek}).
A single \texttt{\textbackslash{}} is a special character in R.
\end{alertblock}
\end{frame}

\begin{frame}[fragile]{The Working Directory}
\protect\hypertarget{the-working-directory}{}
The folder R treats as ``here.'' Every relative path starts from it.

\begin{Shaded}
\begin{Highlighting}[]
\FunctionTok{getwd}\NormalTok{()}
\end{Highlighting}
\end{Shaded}

\begin{verbatim}
[1] "C:/Users/USER/Desktop/DBA_R_Workshop/working"
\end{verbatim}

\begin{Shaded}
\begin{Highlighting}[]
\FunctionTok{setwd}\NormalTok{(}\StringTok{"C:/Users/tek/Documents/dba"}\NormalTok{)   }\CommentTok{\# Windows}
\FunctionTok{setwd}\NormalTok{(}\StringTok{"/Users/tek/Documents/dba"}\NormalTok{)     }\CommentTok{\# Mac}
\FunctionTok{setwd}\NormalTok{(}\StringTok{".."}\NormalTok{)                           }\CommentTok{\# up one folder}
\end{Highlighting}
\end{Shaded}

Or: \textbf{Session \(\rightarrow\) Set Working Directory
\(\rightarrow\) Choose Directory}
\end{frame}

\begin{frame}[fragile]{Why \texttt{setwd()} Causes Problems}
\protect\hypertarget{why-setwd-causes-problems}{}
A script starting with \texttt{setwd("C:/Users/tek/...")} only works on
\textbf{Tek's computer}.

It breaks when you:

\begin{itemize}
\tightlist
\item
  Share it with a coauthor
\item
  Move to a new laptop or rename a folder
\item
  Switch between Windows and Mac
\end{itemize}

\begin{exampleblock}{The solution}
\protect\hypertarget{the-solution}{}
Never type your computer's address into a script. Use \textbf{RStudio
Projects}.
\end{exampleblock}
\end{frame}

\begin{frame}[fragile]{RStudio Projects}
\protect\hypertarget{sec-projects}{}
\begin{itemize}
\tightlist
\item
  A folder with a small \texttt{.Rproj} file inside
\item
  Opening the project sets the working directory to that folder
  \textbf{automatically}
\item
  All paths can be relative, so the folder can be moved, shared, or
  zipped
\end{itemize}

\textbf{Create one:} File \(\rightarrow\) New Project \(\rightarrow\)
New Directory \(\rightarrow\) New Project \(\rightarrow\) name it
\texttt{dba-r-workshop} \(\rightarrow\) Create Project

\textbf{Reopen it:} double-click the \texttt{.Rproj} file, or File
\(\rightarrow\) Recent Projects
\end{frame}

\begin{frame}[fragile]{A Folder Layout That Works}
\protect\hypertarget{a-folder-layout-that-works}{}
The same kind of tree, with the project folder at the top. Paths inside
it stay the same wherever the folder sits.

\begin{verbatim}
dba-r-workshop/
|-- dba-r-workshop.Rproj
|-- data/          # raw data: never edit by hand
|-- output/        # tables, figures, cleaned data your code creates
|-- session1.qmd
`-- session2.qmd
\end{verbatim}

\begin{alertblock}{Avoid OneDrive, iCloud, and Dropbox}
\protect\hypertarget{avoid-onedrive-icloud-and-dropbox}{}
Synced folders can lock files and cause confusing errors. On Windows,
\textbf{Documents} is often synced to OneDrive. Use
\texttt{C:/R-projects} (Windows) or
\texttt{\textasciitilde{}/R-projects} (Mac).
\end{alertblock}
\end{frame}

\begin{frame}[fragile]{The \texttt{here} Package}
\protect\hypertarget{sec-here}{}
\begin{itemize}
\tightlist
\item
  A Quarto document uses \textbf{its own folder} as the working
  directory
\item
  \texttt{here()} always builds paths from the \textbf{project root},
  wherever your file is
\end{itemize}

\begin{Shaded}
\begin{Highlighting}[]
\FunctionTok{library}\NormalTok{(here)}
\end{Highlighting}
\end{Shaded}

\begin{verbatim}
here() starts at C:/Users/USER/Desktop/DBA_R_Workshop
\end{verbatim}

\begin{Shaded}
\begin{Highlighting}[]
\FunctionTok{here}\NormalTok{()                          }\CommentTok{\# the project root}
\end{Highlighting}
\end{Shaded}

\begin{verbatim}
[1] "C:/Users/USER/Desktop/DBA_R_Workshop"
\end{verbatim}

\begin{Shaded}
\begin{Highlighting}[]
\FunctionTok{here}\NormalTok{(}\StringTok{"data"}\NormalTok{, }\StringTok{"survey.csv"}\NormalTok{)      }\CommentTok{\# correct separators on any computer}
\end{Highlighting}
\end{Shaded}

\begin{verbatim}
[1] "C:/Users/USER/Desktop/DBA_R_Workshop/data/survey.csv"
\end{verbatim}

\begin{exampleblock}{The habit to build}
\protect\hypertarget{the-habit-to-build}{}
One Project per research project · raw data in \texttt{data/} · paths
with \texttt{here()} · no \texttt{setwd()} in scripts you keep
\end{exampleblock}
\end{frame}

\hypertarget{sec-files}{%
\section{Accessing Files from R}\label{sec-files}}

\begin{frame}[fragile]{Listing and Creating Files}
\protect\hypertarget{listing-and-creating-files}{}
\begin{Shaded}
\begin{Highlighting}[]
\FunctionTok{list.files}\NormalTok{()                         }\CommentTok{\# files in the working directory}
\end{Highlighting}
\end{Shaded}

\begin{verbatim}
 [1] "session1-slides.aux"       "session1-slides.log"      
 [3] "session1-slides.pdf"       "session1-slides.qmd"      
 [5] "session1-slides.rmarkdown" "session1-slides.tex"      
 [7] "session1-slides.vrb"       "session1-slides_files"    
 [9] "session1.html"             "session1.qmd"             
[11] "session2-slides.qmd"       "session2.qmd"             
\end{verbatim}

\begin{Shaded}
\begin{Highlighting}[]
\FunctionTok{list.files}\NormalTok{(}\AttributeTok{pattern =} \StringTok{"}\SpecialCharTok{\textbackslash{}\textbackslash{}}\StringTok{.qmd$"}\NormalTok{)      }\CommentTok{\# only files ending in .qmd}
\end{Highlighting}
\end{Shaded}

\begin{verbatim}
[1] "session1-slides.qmd" "session1.qmd"        "session2-slides.qmd"
[4] "session2.qmd"       
\end{verbatim}

\begin{Shaded}
\begin{Highlighting}[]
\FunctionTok{dir.create}\NormalTok{(}\FunctionTok{here}\NormalTok{(}\StringTok{"data"}\NormalTok{), }\AttributeTok{showWarnings =} \ConstantTok{FALSE}\NormalTok{)   }\CommentTok{\# no warning if it exists}
\FunctionTok{dir.exists}\NormalTok{(}\FunctionTok{here}\NormalTok{(}\StringTok{"data"}\NormalTok{))}
\end{Highlighting}
\end{Shaded}

\begin{verbatim}
[1] TRUE
\end{verbatim}
\end{frame}

\begin{frame}[fragile]{Writing a CSV File}
\protect\hypertarget{writing-a-csv-file}{}
\begin{Shaded}
\begin{Highlighting}[]
\NormalTok{sales }\OtherTok{\textless{}{-}} \FunctionTok{data.frame}\NormalTok{(}
  \AttributeTok{region  =} \FunctionTok{c}\NormalTok{(}\StringTok{"East"}\NormalTok{, }\StringTok{"South"}\NormalTok{, }\StringTok{"West"}\NormalTok{, }\StringTok{"North"}\NormalTok{),}
  \AttributeTok{revenue =} \FunctionTok{c}\NormalTok{(}\DecValTok{125000}\NormalTok{, }\DecValTok{98000}\NormalTok{, }\DecValTok{143500}\NormalTok{, }\DecValTok{87250}\NormalTok{),}
  \AttributeTok{stores  =} \FunctionTok{c}\NormalTok{(}\DecValTok{12}\NormalTok{, }\DecValTok{9}\NormalTok{, }\DecValTok{15}\NormalTok{, }\DecValTok{8}\NormalTok{)}
\NormalTok{)}

\FunctionTok{write.csv}\NormalTok{(sales, }\FunctionTok{here}\NormalTok{(}\StringTok{"data"}\NormalTok{, }\StringTok{"sales.csv"}\NormalTok{), }\AttributeTok{row.names =} \ConstantTok{FALSE}\NormalTok{)}

\FunctionTok{file.exists}\NormalTok{(}\FunctionTok{here}\NormalTok{(}\StringTok{"data"}\NormalTok{, }\StringTok{"sales.csv"}\NormalTok{))}
\end{Highlighting}
\end{Shaded}

\begin{verbatim}
[1] TRUE
\end{verbatim}

\begin{itemize}
\tightlist
\item
  Open the \texttt{data} folder in the \textbf{Files} pane:
  \texttt{sales.csv} is there
\end{itemize}
\end{frame}

\begin{frame}[fragile]{Reading a CSV File}
\protect\hypertarget{reading-a-csv-file}{}
\begin{Shaded}
\begin{Highlighting}[]
\NormalTok{sales\_in }\OtherTok{\textless{}{-}} \FunctionTok{read.csv}\NormalTok{(}\FunctionTok{here}\NormalTok{(}\StringTok{"data"}\NormalTok{, }\StringTok{"sales.csv"}\NormalTok{))}
\FunctionTok{str}\NormalTok{(sales\_in)        }\CommentTok{\# structure: column names and types}
\end{Highlighting}
\end{Shaded}

\begin{verbatim}
'data.frame':   4 obs. of  3 variables:
 $ region : chr  "East" "South" "West" "North"
 $ revenue: int  125000 98000 143500 87250
 $ stores : int  12 9 15 8
\end{verbatim}

\begin{Shaded}
\begin{Highlighting}[]
\FunctionTok{dim}\NormalTok{(sales\_in)        }\CommentTok{\# rows and columns}
\end{Highlighting}
\end{Shaded}

\begin{verbatim}
[1] 4 3
\end{verbatim}

\begin{itemize}
\tightlist
\item
  Also: \texttt{head()} for the first rows, \texttt{View()} for a
  spreadsheet viewer
\item
  Session 4 covers Excel, SPSS, Stata, and messier CSV files
\end{itemize}

\begin{alertblock}{Deleting files}
\protect\hypertarget{deleting-files}{}
\texttt{file.remove()} and \texttt{unlink()} delete
\textbf{permanently}, without the Recycle Bin or Trash.
\end{alertblock}
\end{frame}

\hypertarget{sec-quarto}{%
\section{Scripts and Quarto Documents}\label{sec-quarto}}

\begin{frame}[fragile]{Two Kinds of Files}
\protect\hypertarget{two-kinds-of-files}{}
\begin{longtable}[]{@{}
  >{\raggedright\arraybackslash}p{(\columnwidth - 4\tabcolsep) * \real{0.3333}}
  >{\raggedright\arraybackslash}p{(\columnwidth - 4\tabcolsep) * \real{0.3333}}
  >{\raggedright\arraybackslash}p{(\columnwidth - 4\tabcolsep) * \real{0.3333}}@{}}
\toprule\noalign{}
\begin{minipage}[b]{\linewidth}\raggedright
\end{minipage} & \begin{minipage}[b]{\linewidth}\raggedright
R script (\texttt{.R})
\end{minipage} & \begin{minipage}[b]{\linewidth}\raggedright
Quarto document (\texttt{.qmd})
\end{minipage} \\
\midrule\noalign{}
\endhead
Contains & Code and \texttt{\#} comments & Text, code, and output
together \\
Best for & Analysis pipelines, data cleaning & Reports, notes, papers,
homework \\
Output & Results in the Console & HTML, Word, or PDF \\
\bottomrule\noalign{}
\end{longtable}

\begin{itemize}
\tightlist
\item
  Quarto is the successor to R Markdown: same ideas
\end{itemize}

\textbf{Your first document:} File \(\rightarrow\) New File
\(\rightarrow\) Quarto Document \(\rightarrow\) title \(\rightarrow\)
\textbf{Create} \(\rightarrow\) save in your project \(\rightarrow\)
\textbf{Render} (\texttt{Ctrl}/\texttt{Cmd} + \texttt{Shift} +
\texttt{K})
\end{frame}

\begin{frame}[fragile]{Anatomy of a Quarto Document}
\protect\hypertarget{anatomy-of-a-quarto-document}{}
\begin{Shaded}
\begin{Highlighting}[]
\CommentTok{{-}{-}{-}}
\AnnotationTok{title:}\CommentTok{ "My First Quarto Document"}
\AnnotationTok{author:}\CommentTok{ "Your Name"}
\AnnotationTok{format:}\CommentTok{ html}
\CommentTok{{-}{-}{-}}

\FunctionTok{\#\# Introduction}

\NormalTok{This is **plain text** written in Markdown.}

\InformationTok{\textasciigrave{}\textasciigrave{}\textasciigrave{}\{r\}}
\CommentTok{\#| label: first{-}chunk}

\NormalTok{x }\OtherTok{\textless{}{-}} \FunctionTok{c}\NormalTok{(}\DecValTok{3}\NormalTok{, }\DecValTok{7}\NormalTok{, }\DecValTok{12}\NormalTok{)}
\FunctionTok{mean}\NormalTok{(x)}
\InformationTok{\textasciigrave{}\textasciigrave{}\textasciigrave{}}
\end{Highlighting}
\end{Shaded}

\begin{enumerate}
\tightlist
\item
  \textbf{Header} between \texttt{-\/-\/-} lines: title and output
  format
\item
  \textbf{Text} in Markdown: \texttt{**bold**}, \texttt{*italic*},
  \texttt{\#\#\ Heading}, \texttt{-\ bullet}
\item
  \textbf{Code chunks}; lines starting with \texttt{\#\textbar{}} are
  chunk options
\end{enumerate}
\end{frame}

\begin{frame}[fragile]{Chunk Options}
\protect\hypertarget{chunk-options}{}
\begin{longtable}[]{@{}ll@{}}
\toprule\noalign{}
Option & Effect \\
\midrule\noalign{}
\endhead
\texttt{\#\textbar{}\ eval:\ false} & Show the code but do not run it \\
\texttt{\#\textbar{}\ echo:\ false} & Run the code, show only the
output \\
\texttt{\#\textbar{}\ include:\ false} & Run the code, show nothing \\
\texttt{\#\textbar{}\ warning:\ false} & Hide warnings \\
\texttt{\#\textbar{}\ label:\ name} & Name the chunk \\
\bottomrule\noalign{}
\end{longtable}

Insert a chunk: \texttt{Ctrl} + \texttt{Alt} + \texttt{I} (Windows),
\texttt{Cmd} + \texttt{Option} + \texttt{I} (Mac)

\begin{alertblock}{Rendering starts from a clean slate}
\protect\hypertarget{rendering-starts-from-a-clean-slate}{}
Render runs your document top to bottom in a \textbf{fresh R session}.
Objects created in the Console do not exist there. ``Object not found''
means the code that creates it is missing from the document.
\end{alertblock}
\end{frame}

\hypertarget{sec-help}{%
\section{Getting Help}\label{sec-help}}

\begin{frame}[fragile]{Help Inside R}
\protect\hypertarget{help-inside-r}{}
\begin{Shaded}
\begin{Highlighting}[]
\NormalTok{?mean                         }\CommentTok{\# help page for a function}
\FunctionTok{help}\NormalTok{(}\AttributeTok{package =} \StringTok{"here"}\NormalTok{)        }\CommentTok{\# everything in a package}
\NormalTok{??}\StringTok{"regression"}                \CommentTok{\# search all installed help pages}
\FunctionTok{vignette}\NormalTok{(}\AttributeTok{package =} \StringTok{"dplyr"}\NormalTok{)   }\CommentTok{\# list a package\textquotesingle{}s tutorials}
\FunctionTok{vignette}\NormalTok{(}\StringTok{"dplyr"}\NormalTok{)             }\CommentTok{\# open one}
\FunctionTok{example}\NormalTok{(mean)                 }\CommentTok{\# run the examples from a help page}
\end{Highlighting}
\end{Shaded}

\textbf{Reading a help page:} Description \(\rightarrow\) Usage
\(\rightarrow\) Arguments \(\rightarrow\) \textbf{Examples} (copy and
run them)
\end{frame}

\begin{frame}[fragile]{Reading Error Messages}
\protect\hypertarget{reading-error-messages}{}
\begin{longtable}[]{@{}
  >{\raggedright\arraybackslash}p{(\columnwidth - 2\tabcolsep) * \real{0.4091}}
  >{\raggedright\arraybackslash}p{(\columnwidth - 2\tabcolsep) * \real{0.5909}}@{}}
\toprule\noalign{}
\begin{minipage}[b]{\linewidth}\raggedright
Message
\end{minipage} & \begin{minipage}[b]{\linewidth}\raggedright
Usual cause
\end{minipage} \\
\midrule\noalign{}
\endhead
\texttt{could\ not\ find\ function\ "ggplot"} & Package not loaded:
\texttt{library(ggplot2)} \\
\texttt{there\ is\ no\ package\ called\ \textquotesingle{}here\textquotesingle{}}
& Not installed: \texttt{install.packages("here")} \\
\texttt{object\ \textquotesingle{}sales\textquotesingle{}\ not\ found} &
Typo, or the code that creates it was not run \\
\texttt{cannot\ open\ file\ ...\ No\ such\ file\ or\ directory} & Wrong
path: check \texttt{getwd()}, use \texttt{here()} \\
\texttt{unexpected\ symbol} /
\texttt{unexpected\ \textquotesingle{})\textquotesingle{}} & Missing
comma, quote, or bracket \\
\bottomrule\noalign{}
\end{longtable}

\begin{exampleblock}{Warnings are not errors}
\protect\hypertarget{warnings-are-not-errors}{}
A \textbf{warning} means R did something but wants you to check it. An
\textbf{error} means R stopped.
\end{exampleblock}
\end{frame}

\begin{frame}[fragile]{Asking for Help Well}
\protect\hypertarget{asking-for-help-well}{}
Include:

\begin{enumerate}
\tightlist
\item
  \textbf{What you were trying to do}
\item
  \textbf{The code}, as text rather than a screenshot
\item
  \textbf{The exact error message}, copied in full
\item
  \textbf{\texttt{sessionInfo()}} output
\end{enumerate}

Writing the question carefully often helps you find the answer yourself.
\end{frame}

\begin{frame}{Using AI Assistants}
\protect\hypertarget{using-ai-assistants}{}
\begin{itemize}
\tightlist
\item
  Very good at explaining error messages and suggesting code
\item
  Treat answers like advice from a helpful classmate:

  \begin{itemize}
  \tightlist
  \item
    Run the code and check the results make sense
  \item
    Understand each line before it goes into your research
  \item
    They can invent functions or arguments that do not exist
  \end{itemize}
\end{itemize}

\begin{alertblock}{Research data}
\protect\hypertarget{research-data}{}
Never paste confidential or identifiable research data into an AI tool
without checking your IRB protocol and your institution's policy.
\end{alertblock}
\end{frame}

\begin{frame}{Good Places to Learn More}
\protect\hypertarget{good-places-to-learn-more}{}
\begin{itemize}
\tightlist
\item
  \href{https://r4ds.hadley.nz/}{R for Data Science (2e)}: free online
  book, the best next step
\item
  \href{https://posit.co/resources/cheatsheets/}{Posit Cheatsheets}:
  one-page summaries
\item
  \href{https://quarto.org/docs/guide/}{Quarto Guide}
\item
  \href{https://forum.posit.co/}{Posit Community} and
  \href{https://stackoverflow.com/questions/tagged/r}{Stack Overflow}:
  search before posting
\end{itemize}
\end{frame}

\hypertarget{sec-troubleshoot}{%
\section{Troubleshooting Installation}\label{sec-troubleshoot}}

\begin{frame}[fragile]{Common Problems}
\protect\hypertarget{common-problems}{}
\begin{longtable}[]{@{}
  >{\raggedright\arraybackslash}p{(\columnwidth - 2\tabcolsep) * \real{0.4091}}
  >{\raggedright\arraybackslash}p{(\columnwidth - 2\tabcolsep) * \real{0.5909}}@{}}
\toprule\noalign{}
\begin{minipage}[b]{\linewidth}\raggedright
Problem
\end{minipage} & \begin{minipage}[b]{\linewidth}\raggedright
What to try
\end{minipage} \\
\midrule\noalign{}
\endhead
Installer asks for an administrator password & Ask IT, or use Posit
Cloud \\
``Use a personal library instead?'' & Answer \textbf{yes} \\
Package errors mentioning OneDrive & Keep projects outside OneDrive \\
Errors with strange characters in paths & Use a simple folder such as
\texttt{C:/R-projects} \\
Mac: ``incompatible architecture'' & Check \texttt{R.version\$arch};
reinstall \texttt{arm64} \\
Mac: arm64 installer refuses & Update to macOS 14+, or Posit Cloud \\
RStudio cannot find R & Install R, then restart the computer \\
A package needs ``compilation'' and fails & Answer \textbf{no}, or
install Rtools / Command Line Tools \\
\bottomrule\noalign{}
\end{longtable}

Where your packages are stored: \texttt{.libPaths()}
\end{frame}

\hypertarget{sec-your-turn}{%
\section{Your Turn}\label{sec-your-turn}}

\begin{frame}[fragile]{Exercise}
\protect\hypertarget{exercise}{}
\begin{alertblock}{Tasks}
\protect\hypertarget{tasks}{}
\begin{enumerate}
\tightlist
\item
  Install R and RStudio. Confirm \texttt{R.version\$arch} matches your
  computer
\item
  Change the RStudio day-one settings
\item
  Create an RStudio Project \texttt{dba-r-workshop}, \textbf{not} in
  OneDrive or iCloud
\item
  Install the course packages; confirm every package returns
  \texttt{TRUE}
\item
  Create a \texttt{data} folder; save the data frame below as
  \texttt{data/employees.csv}; read it back and check it with
  \texttt{str()}
\item
  Create a Quarto document that reads the file and computes the average
  salary; render it
\item
  \textbf{Bonus:} run \texttt{library(heer)} and read the error. Use
  \texttt{?read.csv} to find what \texttt{na.strings} does
\end{enumerate}
\end{alertblock}
\end{frame}

\begin{frame}[fragile]{Exercise Data}
\protect\hypertarget{exercise-data}{}
\begin{Shaded}
\begin{Highlighting}[]
\NormalTok{employees }\OtherTok{\textless{}{-}} \FunctionTok{data.frame}\NormalTok{(}
  \AttributeTok{name       =} \FunctionTok{c}\NormalTok{(}\StringTok{"Ana"}\NormalTok{, }\StringTok{"Ben"}\NormalTok{, }\StringTok{"Chloe"}\NormalTok{, }\StringTok{"Dev"}\NormalTok{, }\StringTok{"Eun{-}ji"}\NormalTok{),}
  \AttributeTok{department =} \FunctionTok{c}\NormalTok{(}\StringTok{"Sales"}\NormalTok{, }\StringTok{"HR"}\NormalTok{, }\StringTok{"IT"}\NormalTok{, }\StringTok{"Sales"}\NormalTok{, }\StringTok{"IT"}\NormalTok{),}
  \AttributeTok{salary     =} \FunctionTok{c}\NormalTok{(}\DecValTok{68000}\NormalTok{, }\DecValTok{72000}\NormalTok{, }\DecValTok{95000}\NormalTok{, }\DecValTok{61000}\NormalTok{, }\DecValTok{102000}\NormalTok{)}
\NormalTok{)}
\end{Highlighting}
\end{Shaded}
\end{frame}

\begin{frame}[fragile]{Exercise: Solution (Steps 5 and 6)}
\protect\hypertarget{exercise-solution-steps-5-and-6}{}
\begin{Shaded}
\begin{Highlighting}[]
\FunctionTok{library}\NormalTok{(here)}

\FunctionTok{dir.create}\NormalTok{(}\FunctionTok{here}\NormalTok{(}\StringTok{"data"}\NormalTok{), }\AttributeTok{showWarnings =} \ConstantTok{FALSE}\NormalTok{)}
\FunctionTok{write.csv}\NormalTok{(employees, }\FunctionTok{here}\NormalTok{(}\StringTok{"data"}\NormalTok{, }\StringTok{"employees.csv"}\NormalTok{), }\AttributeTok{row.names =} \ConstantTok{FALSE}\NormalTok{)}

\NormalTok{employees\_in }\OtherTok{\textless{}{-}} \FunctionTok{read.csv}\NormalTok{(}\FunctionTok{here}\NormalTok{(}\StringTok{"data"}\NormalTok{, }\StringTok{"employees.csv"}\NormalTok{))}
\FunctionTok{str}\NormalTok{(employees\_in)}
\end{Highlighting}
\end{Shaded}

\begin{verbatim}
'data.frame':   5 obs. of  3 variables:
 $ name      : chr  "Ana" "Ben" "Chloe" "Dev" ...
 $ department: chr  "Sales" "HR" "IT" "Sales" ...
 $ salary    : int  68000 72000 95000 61000 102000
\end{verbatim}

\begin{Shaded}
\begin{Highlighting}[]
\FunctionTok{mean}\NormalTok{(employees\_in}\SpecialCharTok{$}\NormalTok{salary)}
\end{Highlighting}
\end{Shaded}

\begin{verbatim}
[1] 79600
\end{verbatim}

\begin{itemize}
\tightlist
\item
  The Quarto document needs \textbf{all} of this code, including
  \texttt{library(here)}, because rendering starts from a fresh R
  session
\end{itemize}
\end{frame}

\hypertarget{sec-summary}{%
\section{Summary}\label{sec-summary}}

\begin{frame}[fragile]{Key Takeaways}
\protect\hypertarget{key-takeaways}{}
\begin{itemize}
\tightlist
\item
  \textbf{R is the engine, RStudio is the dashboard.} Install R for your
  processor; on Apple silicon, \texttt{R.version\$arch} must say
  \texttt{"aarch64"}
\item
  \textbf{Turn off saving and restoring \texttt{.RData}}: scripts are
  the complete record of your work
\item
  \textbf{Install} once with \texttt{install.packages()}; \textbf{load}
  every session with \texttt{library()}
\item
  \textbf{One RStudio Project per research project}, paths with
  \texttt{here()}, no \texttt{setwd()} in scripts you keep
\item
  \textbf{Quarto renders in a fresh R session}: documents must contain
  all the code they need
\item
  \textbf{Read error messages}, use \texttt{?function}, and include
  code, error, and \texttt{sessionInfo()} when asking for help
\end{itemize}

\textbf{Next session:} data types and structures
\end{frame}



\end{document}
